% Part: first-order-logic
% Chapter: model-theory
% Section: overspill

\documentclass[../../../include/open-logic-section]{subfiles}

\begin{document}

\olfileid{mod}{bas}{ove}
\olsection{ওভারস্পিল}

\begin{thm}
\ollabel{overspill} কোনো বাক্যসমষ্টি~$\Gamma$-র ইচ্ছামতো বড় সসীম মডেল
থাকলে তার একটি অসীম মডেলও আছে।
\end{thm}

\begin{proof}
$\Gamma$-র ভাষায় গণনাযোগ্য অসীমসংখ্যক নতুন ধ্রুবক $c_0$, $c_1$,
\dots{} যোগ করে ভাষাটি সম্প্রসারিত করো এবং সেট $\Gamma \cup \{c_i
\neq c_j : i \neq j\}$ বিবেচনা করো। $\Gamma$-র ইচ্ছামতো বড় সসীম
মডেল আছে বলার অর্থ হলো, প্রত্যেক $m >0$-এর জন্য এমন $n\ge m$ আছে যেন
$\Gamma$-র $n$ আকারের একটি মডেল থাকে। এর ফলে $\Gamma \cup \{c_i
\neq c_j : i \neq j\}$ সসীমভাবে পরিতৃপ্তিযোগ্য। সংহতি অনুসারে,
$\Gamma \cup \{c_i \neq c_j : i \neq j\}$-এর একটি মডেল~$\Struct M$
আছে; তার সংজ্ঞাক্ষেত্র অবশ্যই অসীম, কারণ সেটি সব অসমতা $c_i \neq c_j$
পরিতৃপ্ত করে।
\end{proof}

\begin{prop}
\ollabel{inf-not-fo}
কোনো প্রথম-ক্রমের ভাষাতেই এমন কোনো বাক্য~$!A$ নেই, যা কোনো
!!a{structure}~$\Struct M$-এ ঠিক তখনই সত্য, যখন !!{structure}-টির
সংজ্ঞাক্ষেত্র~$\Domain{M}$ অসীম।
\end{prop}

\begin{proof}
এমন কোনো~$!A$ থাকলে তার নঞর্থকরণ~$\lnot !A$ ঠিক সব সসীম
!!{structure}s-এ সত্য হতো। ফলে তার ইচ্ছামতো বড় সসীম মডেল থাকত, কিন্তু
কোনো অসীম মডেল থাকত না; এটি \olref{overspill}-এর বিরোধী।
\end{proof}

\end{document}
